\metatitle{Ehrhart theory}
\metadescription{Ehrhart polynomials, Ehrhart series, h-star vectors,
reciprocity, and common appearances in algebraic combinatorics.}
\metakeywords{Ehrhart theory,Ehrhart polynomial,Ehrhart series,h-star vector,
Ehrhart--Macdonald reciprocity,order polytope}

\section[ehrhart]{Ehrhart theory}

\defin{Ehrhart theory}\label{ehrhartTheory} studies the lattice points in
dilates of rational and lattice polytopes.  For a lattice polytope \(P\), the
basic counting function is
\[
  E_P(n)=|nP\cap \setZ^d|.
\]
This is the \hyperref[ehrhartPolynomialDefinition]{Ehrhart polynomial} of
\(P\).  Its generating function is the
\hyperref[ehrhartSeries]{Ehrhart series}, usually written in the form
\[
  \sum_{n\geq 0} E_P(n)t^n
  =
  \frac{h^*_P(t)}{(1-t)^{d+1}},
\]
where \(h^*_P(t)\) is the
\hyperref[hStarPolynomial]{\(h^*\)-polynomial}.  See
\cite{Stanley1993,Beck2007} and the \hyperref[ehrhartPolytopes]{polytope}
page for the standard definitions and basic theorems.

\begin{example}
For the unit interval \(P=[0,1]\),
\[
  E_P(n)=n+1,\qquad
  \sum_{n\geq 0}(n+1)t^n=\frac{1}{(1-t)^2}.
\]
Thus \(h^*_P(t)=1\).  For the unit square \([0,1]^2\),
\[
  E_P(n)=(n+1)^2,\qquad
  \sum_{n\geq 0}(n+1)^2t^n=\frac{1+t}{(1-t)^3},
\]
so \(h^*_P(t)=1+t\).
\end{example}

\begin{example}
For the \(d\)-dimensional cube \([0,1]^d\), the \(h^*\)-polynomial is the
Eulerian polynomial:
\[
\begin{array}{c|l}
d & h^*_{[0,1]^d}(t)\\
\hline
1 & 1\\
2 & 1+t\\
3 & 1+4t+t^2\\
4 & 1+11t+11t^2+t^3.
\end{array}
\]
The coefficients form the Eulerian number triangle \oeis{A008292}.  This is
one of the quickest ways Ehrhart theory meets permutation statistics.
\end{example}

\subsection[ehrhartReciprocityHub]{Reciprocity and duality}

\hyperref[ehrhartMacdonaldReciprocity]{Ehrhart--Macdonald reciprocity} relates
lattice points in \(nP\) to lattice points in the interior of \(nP\), by
evaluating the Ehrhart polynomial at negative integers.  For a
\(d\)-dimensional lattice polytope,
\[
  E_P(-n)=(-1)^d |nP^\circ\cap \setZ^d|.
\]
In terms of the Ehrhart series, palindromicity of the
\hyperref[hStarVector]{\(h^*\)-vector} is tied to the
\hyperref[gorenstein]{Gorenstein} condition.  This is one reason Ehrhart
theory often behaves like a discrete shadow of Poincaré duality.

\subsection[hStarInequalities]{Inequalities for \(h^*\)-vectors}

The \hyperref[hStarVector]{\(h^*\)-vector}
\[
  h^*(P)=(h^*_0,h^*_1,\dotsc,h^*_d)
\]
of a \(d\)-dimensional lattice polytope is not an arbitrary nonnegative
integer vector.  The first restrictions are the basic ones:
\[
  h^*_0=1,\qquad
  h^*_1=|P\cap \setZ^d|-d-1,\qquad
  \sum_{i=0}^{d} h_i^*=\operatorname{Vol}(P),
\]
where \(\operatorname{Vol}(P)\) is the normalized volume.  If \(s\) is the
degree of \(h^*_P(t)\), then \(h^*_s\) counts lattice points in a suitable
interior dilation of \(P\), and in particular
\[
  h^*_d=|P^\circ\cap \setZ^d|.
\]
This gives the elementary inequality
\[
  h^*_d\leq h^*_1.
\]

There are also less immediate inequalities.  \name{Takayuki Hibi} proved
\cite{Hibi1994LowerBound} that
\[
  h^*_{d-1}+h^*_{d-2}+\dotsb+h^*_{d-i}
  \leq
  h^*_2+h^*_3+\dotsb+h^*_{i+1},
  \qquad 1\leq i\leq d-1.
\]
If \(P\) has an interior lattice point, equivalently \(h^*_d>0\), then Hibi
also proved
\[
  h^*_1\leq h^*_i,\qquad 1\leq i\leq d-1.
\]
Another family, due to \name{Richard P. Stanley}
\cite{Stanley1991HilbertFunction}, says that if \(s=\deg h^*_P(t)\), then
\[
  h^*_0+h^*_1+\dotsb+h^*_i
  \leq
  h^*_s+h^*_{s-1}+\dotsb+h^*_{s-i},
  \qquad 0\leq i\leq s.
\]
These inequalities are useful quick checks when an experimental
\(h^*\)-vector has been computed.

The two-dimensional case has a particularly concrete form.  Scott's theorem
\cite{Scott1976ConvexLatticePolygons} characterizes the possible
\(h^*\)-polynomials \(1+h^*_1t+h^*_2t^2\) of lattice polygons by the
following alternatives:
\[
  h^*_2=0,\qquad
  h^*_2\leq h^*_1\leq 3h^*_2+3,\qquad
  \text{or}\qquad (h^*_1,h^*_2)=(7,1).
\]
\name{Jaron Treutlein} extended the necessary part of this statement to
lattice polytopes of degree at most two \cite{Treutlein2010DegreeTwo}.

Finally, \name{Gabriele Balletti} and \name{Akihiro Higashitani}
\cite{BallettiHigashitani2018Universal} proved a universal version of Scott's
inequality.  If \(P\) is any lattice polytope, in any
dimension and of any degree, and
\[
  h^*_P(t)=1+h^*_1t+h^*_2t^2+\dotsb
  \qquad\text{satisfies}\qquad h^*_3=0,
\]
then one of the following holds:
\[
  h^*_2=0,\qquad
  h^*_1\leq 3h^*_2+3,\qquad
  \text{or}\qquad (h^*_1,h^*_2)=(7,1).
\]
The hypothesis \(h^*_3=0\) is essential: their paper gives a
five-dimensional example with \(h^*_P(t)=1+8t+t^2+8t^3\), which violates
Scott's inequality once the \(h^*_3\)-term is present.  Stapledon's work
\cite{Stapledon2009Inequalities,Stapledon2016Additive} gives further
families of inequalities, often sharper but less uniform in appearance.

\subsection[ehrhartPosets]{Order polytopes and posets}

For a finite \hyperref[poset]{poset} \(P\), the
\hyperref[orderPolytope]{order polytope} \(\mathcal O(P)\) packages
order-preserving maps \(P\to[0,1]\).  Its \(h^*\)-polynomial is the
\hyperref[pEulerianPolynomial]{\(P\)-Eulerian polynomial}; equivalently, it
records descents of linear extensions.  The
\hyperref[chainPolytope]{chain polytope} \(\mathcal C(P)\) has the same
Ehrhart polynomial as \(\mathcal O(P)\), but its inequalities are controlled
by chains rather than order relations.

\todo{Add Yakob Kahane's arXiv:2607.11225v1 on fence-poset order
polynomial coefficients, including the Ehrhart coefficient interpretation for
base polytopes of Schubert matroids and the relevant \(h^*\)-polynomial
consequences.}

\begin{example}
If \(P\) is an antichain on \(d\) elements, then
\(\mathcal O(P)=[0,1]^d\).  Hence
\[
  E_{\mathcal O(P)}(n)=(n+1)^d.
\]
Its \(h^*\)-polynomial is the Eulerian polynomial, whose coefficients are
the Eulerian numbers \oeis{A008292}.  Thus even the most basic order
polytopes already produce familiar sequence data.

If \(P\) is a chain on \(d\) elements, then
\(\mathcal O(P)\) is a simplex, and \(E_{\mathcal O(P)}(n)\) counts weakly
increasing sequences
\[
  0\leq a_1\leq a_2\leq \dotsb \leq a_d\leq n.
\]
\end{example}

\begin{example}
Let \(P_\lambda\) be the Ferrers poset of a partition \(\lambda\).  Linear
extensions of \(P_\lambda\) are
\hyperref[standardYoungTableau]{standard Young tableaux} of shape \(\lambda\),
and the \hyperref[hStarPolynomial]{\(h^*\)-polynomial} of
\(\mathcal O(P_\lambda)\) is the descent
generating polynomial over those tableaux.  For \(\lambda=(n,n)\), this gives
the Narayana polynomial; the Narayana triangle is \oeis{A001263}.  This is
the Ehrhart-theoretic form of the example on the
\hyperref[realRootedTableaux]{real-rooted tableaux page}.
\end{example}

\subsection[ehrhartWhereItAppears]{Where it appears}

Ehrhart theory appears in several parts of algebraic combinatorics:
\begin{itemize}
\item \hyperref[gtpolytopes]{Gelfand--Tsetlin polytopes} and marked order
polytopes model representation-theoretic multiplicities.
\item \hyperref[gtKostkaEhrhart]{Skew and flagged GT-polytopes} give
stretched Kostka, flagged skew Schur, and Littlewood--Richardson-type
counting functions.
\item \hyperref[littlewoodRichardsonModels]{Littlewood--Richardson
coefficients} can be studied through lattice points in hive and related
polytopes.
\item \hyperref[matroidBasePolytope]{Matroid base polytopes} and
\hyperref[positroidPolytopes]{positroid polytopes} have Ehrhart theory tied
to matroid invariants.
\item \hyperref[gammaPositivity]{Gamma-positivity} and
\hyperref[realRootedTableaux]{real-rootedness questions} often ask whether
the corresponding \(h^*\)-polynomials have stronger positivity properties.
\item OEIS comparisons are especially useful for \(h^*\)-triangles arising
from order polytopes, flow polytopes, and lecture-hall-type polytopes, since
their first rows often coincide with classical Eulerian, Narayana, Catalan,
or Euler-number triangles.
\item The \hyperref[transferMatrixMethod]{transfer matrix method} can compute
Ehrhart series when lattice points are built from finitely many boundary
states.
\end{itemize}

\subsection[ehrhartConjecturalPositivity]{Conjectural Ehrhart positivity}

A recurring phenomenon is that representation-theoretic GT-polytopes seem to
have more positive Ehrhart polynomials than arbitrary polytopes.  This should
be treated as a conjectural pattern, since Ehrhart positivity is false for
general lattice polytopes and even for many natural-looking families.

\begin{conjecture}
The Ehrhart polynomial of the skew Gelfand--Tsetlin polytope
\(\gtp_{\lambda/\mu}\) has nonnegative coefficients.  Equivalently,
\[
  k\mapsto \schurS_{k\lambda/k\mu}(1^m)
\]
has nonnegative coefficients as a polynomial in \(k\).  This is the
\hyperref[gtKostkaEhrhart]{Alexandersson--Alhajjar conjecture} on the
Gelfand--Tsetlin page \cite{AlexanderssonAlhajjar2018}.
\end{conjecture}

\begin{conjecture}
The flagged skew Schur stretching polynomial
\[
  k \mapsto K_{k\lambda/k\mu,k\nu}(\avec,\bvec)
\]
has nonnegative coefficients.  This is the
\hyperref[schurFlagged]{flagged generalization} of the
King--Tollu--Toumazet positivity conjecture
\cite{AlexanderssonOguz2023x}.
\end{conjecture}

\begin{conjecture}
For fixed \(\lambda\) and \(\sigma\in\symS_n\), the polynomial
\[
  k\mapsto \key_{k\lambda,\sigma}(1^n)
\]
is Ehrhart-positive; more precisely, the polynomial
\(P_\sigma(\lambda_1,\dotsc,\lambda_n;k)\) from the
\hyperref[keyEhrhart]{GT-face model for key polynomials} has nonnegative
coefficients, and remains nonnegative after changing variables to partial
sums \(a_j=\lambda_1+\dotsb+\lambda_j\)
\cite{AlexanderssonAlhajjar2018}.
\end{conjecture}

These conjectures are parallel: ordinary skew Schur functions, flagged skew
Schur functions, and key polynomials all have GT-type models, and stretching
the highest weights turns the relevant specialization into an Ehrhart
polynomial or a closely related lattice-point count.  The evidence suggests
that the representation-theoretic inequalities impose positivity that is not
visible for arbitrary polytopes.

\subsection[ehrhartQuestions]{Typical questions}

Common questions include:
\begin{itemize}
\item whether the coefficients of \(E_P(n)\) are nonnegative;
\item whether \(h^*_P(t)\) is unimodal, real-rooted, or
\hyperref[gammaPositivePolynomial]{gamma-positive};
\item whether a rational polytope has
\hyperref[periodCollapse]{period collapse};
\item whether a family admits a combinatorial interpretation for its
\(h^*\)-coefficients.
\end{itemize}

For experimental families, the first rows of \(h^*\)-vectors are often among
the best invariants to compare with OEIS.  When an \(h^*\)-triangle is
recognizable, recording the OEIS entry usually helps locate older
enumerative interpretations, recurrences, and transforms.

These questions are often subtle: the \(h^*\)-vector of a lattice polytope is
always nonnegative, but Ehrhart polynomial coefficients can be negative.
